% ===== PRÉAMBULE CANONIQUE — à coller VERBATIM en tête de CHAQUE .tex =====
\documentclass[11pt,a4paper]{article}
\usepackage[utf8]{inputenc}\usepackage[T1]{fontenc}\usepackage{lmodern}
\usepackage[french]{babel}
\usepackage{amsmath,amssymb,mathtools}\usepackage{icomma}
\usepackage{siunitx}\sisetup{locale=FR}  % \qty{}{}, \si{}, \ang{}
\usepackage{xcolor}\usepackage{colortbl}
\usepackage{tikz}\usepackage{pgfplots}\pgfplotsset{compat=1.18}
\usepackage[siunitx,europeanresistors]{circuitikz} % circuits (élec.)
\usepackage[most]{tcolorbox}\usepackage{enumitem}\usepackage{array,booktabs}
\usepackage[a4paper,margin=2.2cm]{geometry}
\usetikzlibrary{arrows.meta,calc,angles,quotes,positioning,patterns,%
  backgrounds,shapes.geometric,decorations.pathmorphing,decorations.markings,intersections}
\definecolor{primary}{RGB}{222,49,99}\definecolor{primarydark}{RGB}{150,22,55}
\definecolor{defcol}{RGB}{0,114,178}\definecolor{methcol}{RGB}{52,92,148}
\definecolor{excol}{RGB}{0,135,90}\definecolor{attncol}{RGB}{200,40,55}
\tcbset{boxbase/.style={enhanced,breakable,boxrule=0.9pt,arc=2.5pt,
  left=3mm,right=3mm,top=2.3mm,bottom=2.3mm,fonttitle=\bfseries,coltitle=white}}
% --- boîtes TUTORIEL (noms EXACTS, ne pas renommer) ---
\newtcolorbox{cle}[1][]{boxbase,colback=defcol!6,colframe=defcol,
  title={\textsc{À retenir}\ifx\relax#1\relax\relax\else\textmd{ : \itshape #1}\fi}}
\newtcolorbox{methode}[1][]{boxbase,colback=methcol!7,colframe=methcol,sharp corners,
  title={\textsc{Méthode}\ifx\relax#1\relax\relax\else\textmd{ --- \itshape #1}\fi}}
\newtcolorbox{exemplebox}[1][]{boxbase,colback=excol!5,colframe=excol,
  title={\textsc{Exemple}\ifx\relax#1\relax\relax\else\textmd{ : \itshape #1}\fi}}
\newtcolorbox{piege}[1][]{boxbase,colback=attncol!6,colframe=attncol,
  title={\textsc{Piège}\ifx\relax#1\relax\relax\else\textmd{ : \itshape #1}\fi}}
\newtcolorbox{remarque}[1][]{boxbase,colback=black!4,colframe=black!45,
  title={\textsc{Remarque}\ifx\relax#1\relax\relax\else\textmd{ : \itshape #1}\fi}}
% --- boîtes EXERCICE / CORRIGÉ (noms EXACTS, auto-comptées) ---
\newtcolorbox[auto counter]{exo}[1][]{boxbase,colback=primary!4,colframe=primary,
  title={\textsc{Exercice}~\thetcbcounter\ifx\relax#1\relax\relax\else\textmd{ --- \itshape #1}\fi}}
\newtcolorbox[auto counter]{corrige}[1][]{boxbase,colback=excol!4,colframe=excol,
  title={\textsc{Corrigé}~\thetcbcounter\ifx\relax#1\relax\relax\else\textmd{ --- \itshape #1}\fi}}
\newcommand{\vv}[1]{\vec{#1}}\newcommand{\ds}{\displaystyle}
\newcommand{\R}{\mathbb{R}}\newcommand{\N}{\mathbb{N}}\newcommand{\Z}{\mathbb{Z}}
\begin{document}
\sisetup{per-mode=symbol}

\begin{exo}[à chaque position son instrument]
Une lentille convergente de focale $f$ est utilisée de trois façons. Associe chaque position de
l'objet à l'instrument correspondant (loupe, appareil photo, projecteur).
\begin{enumerate}[label=\alph*)]
  \item objet au-delà de $2f$ ;
  \item objet entre $f$ et $2f$ ;
  \item objet entre $O$ et $f$.
\end{enumerate}
\end{exo}

\begin{exo}[l'appareil photographique]
\begin{enumerate}[label=\alph*)]
  \item L'image formée sur le capteur est-elle réelle ou virtuelle ? droite ou renversée ?
        agrandie ou réduite ?
  \item Quel composant remplace la pellicule argentique dans un appareil moderne ?
\end{enumerate}
\end{exo}

\begin{exo}[la loupe]
\begin{enumerate}[label=\alph*)]
  \item Par rapport au foyer $F$, où faut-il placer l'objet pour utiliser une lentille convergente en
        loupe ?
  \item Quelle est la nature de l'image (réelle/virtuelle, sens, taille) ?
\end{enumerate}
\end{exo}

\begin{exo}[le microscope optique]
\begin{enumerate}[label=\alph*)]
  \item De combien de lentilles est-il composé, et comment s'appellent-elles ?
  \item Ces lentilles sont-elles convergentes ou divergentes ?
\end{enumerate}
\end{exo}

\begin{exo}[le projecteur]
\begin{enumerate}[label=\alph*)]
  \item Où doit-on placer l'objet (la diapositive) par rapport à $f$ et $2f$ ?
  \item L'image projetée sur l'écran est-elle réelle ou virtuelle ? agrandie ou réduite ?
\end{enumerate}
\end{exo}

\end{document}
